% LTeX: enabled=false

\chapter{Acrònims} %Respectau títol del capítol.
%
% Per utilitzar els acrònims es recomana fer un poc 
% de recerca bibliogràfica per entendre com 
% funcionen. Concretament podeu llegir el manual
% que teniu dins el vostre sistema.
% La comanda `texdoc acronym` hauria de mostrar-lo.
%
\begin{acronym}

\acro{AMC}[AMC]{Adaptive Modulation and Coding}

\acro{EPS}[EPS]{Escola Politècnica Superior}

\acro{IP}[IP]{Internet Protocol}

\acro{LTE}[LTE]{Long Term Evolution}

\acro{MIMO}[MIMO]{Multiple-Input Multiple-Output}

\acro{OFDM}[OFDM]{Orthogonal Frequency Division Multiplexing}

\acro{OFDMA}[OFDMA]{Orthogonal Frequency Division Multiple Access}

\acro{RDI}[R+D+I]{Recerca, Desenvolupament i Innovació}

\acro{TCP}[TCP]{Transport Control Protocol}

\acro{TDMA}[TDMA]{Time Division Multiple Access}

\acro{TFG}[TFG]{Treball Final de Grau}

\acro{GEIN}{Grau d'Enginyeria Informàtica}

\acro{UIB}[UIB]{Universitat de les Illes Balears}

\acro{OUSIS}[OUSIS]{Oficina d'Universitat Saludable i Sostenible}

\acro{LVUIB}[LVUIB]{Liga de Voleibol de la Universitat de les Illes Balears}

\acro{SaaS}[SaaS]{Software as a Service, Software como Servicio}

\acro{OS}[OS]{Object Storage, Almacenamiento de Objetos}

\acro{BD}[BD]{Base de Datos}

\acro{UX}[UX]{User Experience, Experiencia de Usario}

\acro{DTO}[DTO]{Data Transfer Object, Objeto de Transferencia de Datos}

\acro{JVM}[JVM]{Java Virtual Machine, Máquina Virtual de Java}

\acro{ORM}[ORM]{Object-Relational Mapping, Mapeo Relacional de Objetos}

\acro{ABAC}[ABAC]{Attribute-Based Access Control, Control de Acceso Basado en Atributos}

\acro{XSS}[XSS]{Cross-Site Scripting}

\end{acronym}
